\documentclass{article}
\usepackage{tikz}
\usetikzlibrary{angles, quotes}
\usepackage{amssymb}
\usepackage{pgfplots}
\pgfplotsset{width=10cm,compat=1.9}
\usepgfplotslibrary{statistics}
\usepackage[left=2.54cm,right=2.54cm,top=2.54cm,bottom=2.54cm]{geometry}
\usepackage{tasks} %Paquete necesario para la producción de items horizontales para algunos ejercicios de tarea empleando el comando /begin{multicols}}{#}.../end{multicols} para ayudarnos a reducir las páginas
\usepackage{pdflscape} %comando que permite cambiar a orientación horizontal las hojas%
\usepackage{soul} %Paquete para permitir la compilación de acentos usando el comando {\hl{}}}\\
\sethlcolor{yellow(munsell)} %comando para definir el color para el subrayado usando el comando \hl
\usepackage[utf8]{inputenc}
\usepackage[spanish]{babel}
\usepackage{icomma}
\usepackage{ragged2e}
\usepackage{siunitx}
\usepackage{url}
\usepackage[colorlinks=true, urlcolor=blue,  linkcolor=black, citecolor=green]{hyperref}
\usepackage{pdfpages}
\usepackage{blindtext} %Paquete que genera texto ficticio (dummy text) con el comando: \blindtext
\usepackage{cancel} %in the preamble gives you four different modes of striking through: \cancel{text to cancel} draws a diagonal line (slash) through its argument, \bcancel{text to cancel} uses the negative slope (a backslash), \xcancel{text to cancel} draws an X (actually \cancel plus \bcancel), \cancelto{〈value〉}{〈expression〉} draws a diagonal arrow through the 〈expression〉pointing to the 〈value〉 (math-mode only)
\usepackage{csquotes}
\usepackage{afterpage}
\usepackage{parskip} 
\usepackage{float}
\usepackage{enumitem}
\usepackage{multicol}%Paquete que permite la creación aislada de columnas de texto. De esta forma se reduce la cantidad de páginas en nuestro documento
\usepackage{lipsum} %Paquete que genera texto de relleno al igual que \usepackage{blindtext}, se usa con el comando \lipsum[#-#] los #(hashtags) delimitan el número de páginas que deseamos ocupar del paquete lipsum para usarlos como relleno. 

% Definir el comando personalizado para la nota
\newcommand{\mynote}[1]{%
\begin{tikzpicture}[baseline=-0.75ex]
\node[draw, circle, fill=Apple Green!60, inner sep=2pt] (note) {\textbf{Nota:}};
\end{tikzpicture}%
\ \textit{#1}%
}

\newenvironment{Figura}
{\par\medskip\noindent\minipage{\linewidth}}
{\endminipage\par\medskip}
\usepackage{caption}
\usepackage[
backend=biber,
sorting=none,
url=true
]{biblatex} %bibliografía
\addbibresource{biblio.bib}
% Establecer el espacio entre las entradas de la bibliografía
\setlength{\bibitemsep}{1\baselineskip} % Puedes ajustar el valor según tus preferencias
\usepackage{amsmath,amsthm,amssymb,amsfonts}
\usepackage{pifont} %Permite la compilación del comando \xmark: es el símbolo de la tachita
\usepackage{mathtools} %permite la compilación de símbolos para matrices
\usepackage{empheq} %Paquete que se relaciona con \usepackage[most]{tcolorbox} para la creación de cuadros/cajas de colores para delimitar los resultados matemáticos o líneas de texto usando la línea de comando: \begin{empheq}[box={\mymath[colback=orange(webcolor)!70,drop lifted shadow]}]{equation*} \end{empheq}
\usepackage[most]{tcolorbox} %Paquete que se relaciona con \usepackage{empheq} para la creación de cuadros/cajas de colores para delimitar los resultados matemáticos o líneas de texto
\renewcommand{\qedsymbol}{$\blacksquare$}
\usetikzlibrary{positioning,decorations.pathreplacing} %Paquete que permite la compilación de llaves para cuadros sinópticos (brace diagram) 
\usepackage{schemata} %The schemata package is designed just for these brace diagrams
\usepackage{booktabs}

\newcommand\AB[2]{\schema{\schemabox{#1}}{\schemabox{#2}}} %comando o paquete necesario para crear cuadros sinópticos

\newtcbox{\mymath}[1][]{%
nobeforeafter, math upper, tcbox raise base,
enhanced, colframe=blue!30!black,
colback=blue!30, boxrule=1pt,
#1} %Comando importante para encasillar los resultados matemáticos en cajas de diferentes colores, se relaciona con los paquetes \usepackage{empheq} y \usepackage[most]{tcolorbox} 

\newenvironment{sysmatrix}[1]
{\left(\begin{array}{@{}#1@{}}}
{\end{array}\right)}
\newcommand{\ro}[1]{%
\xrightarrow{\mathmakebox[\rowidth]{#1}}%
}
\newlength{\rowidth}% row operation width
\AtBeginDocument{\setlength{\rowidth}{3em}}  %Comando importante para laproducción de lineas de operaciones entre matrices, método de Gauss o Gauss Jordan se relaciona con \newenvironment{sysmatrix}

%formato para cambiar el horario a español
\usepackage[spanish]{datetime2}
\DTMsetdatestyle{spanish}

\renewcommand{\today}{\DTMdisplaydate{\the\year}{\the\month}{\the\day }{-1}}
%%%%%%%%%%%%%%%%%%%%%%%%%%%%%%%%%%%%%%%%%

\usepackage{datetime}
\newcommand{\mycurrenttime}{\xxivtime}

%%%%%%%%%%%%%%%%%%%%%%%%%%%%%%%%%%%%%%%%%%%%%%%%%%%%%%%%%%%%%%%%%%%%%%%%%%%%%%%%%%%%%%%%%%%%%%%%%%%%%%%%%%%%%%%%%%%%%%%%%%%%%%%%%%%%%%%%%%%%%%%%%


%%%%%%%%%%%%%Caja de color para el título%%%%%%%%%%%%%%%%%%%%%%%%%%%%%%%%%%%%%%%%%%%%%%%%%%%%%%

\definecolor{myframecolor}{RGB}{1, 45, 75} % Prussian Blue
\definecolor{myboxcolor}{RGB}{0, 107, 92} % Apple Green

\newtcolorbox{mybox}{
enhanced,
colback=myboxcolor!7, % Color de fondo del cuadro
colframe=myframecolor, % Color del marco del cuadro
arc=0pt,
boxrule=1pt,
borderline west={2mm}{-10mm}{myframecolor}, % Borde en el lado izquierdo
sharp corners=southwest,
width=\linewidth,
before=\par\vspace{\bigskipamount}, % Espacio antes del cuadro
after=\par\vspace{\bigskipamount} % Espacio después del cuadro
}


\newcommand{\euler}{e} %Comando para producir letra e de euler
\newenvironment{remark}{\par\vfill\footnotesize % Vertical white space above the remark and smaller font size
\begin{list}{}{
	\leftmargin=80pt % Indentation on the left
	\rightmargin=60pt}\item\ignorespaces % Indentation on the right
\makebox[-2.5pt]{\begin{tikzpicture}[overlay]
		\node[draw=Horizon!60,line width=2.5pt,circle,fill=Horizon!25,font=\sffamily\bfseries,inner sep=4pt,outer sep=0pt] at (-19pt,5pt){\textcolor{Horizon}{Nota}};\end{tikzpicture}} % Blue Nota in a circle
\advance\baselineskip -1pt}{\end{list}\vskip5pt} % Tighter line spacing and white space after remark
\usepackage{graphicx}

\usepackage{titling}

\input{colores.tex}

%%Producir signos de cita textual``Asignatura''

\newcommand{\xmark}{\ding{55}} %%%comando que genera la tachita

\newenvironment{MyColorPar}[1]{%
\leavevmode\color{#1}\ignorespaces%
}{%
}%

%%%%%%%%%%%%%%%%%%% Título de la tarea, Nombre de alumno

\title{Clase \# 34 - Mecánica Estadística}
\author{\emph{JA} $\boldsymbol{\mid}$}
\date{\today}

%%%%%%%%%%%%%%%%%%% Datos de la Materia

\usepackage{fancyhdr}
\fancypagestyle{plain}{%  the preset of fancyhdr 
\fancyhf{} % clear all header and footer fields
\fancyfoot[R]{
\includegraphics[width=2cm]{LogoFCFMBUAP (1).png}}
\fancyfoot[L]{{\bfseries{\thedate{}}} a las {\bfseries{\mycurrenttime{}}} horas{} (GMT-6, H. Puebla de Zaragoza, Pue)}
\fancyhead[L]{Mecánica Estadística}
\fancyhead[R]{\theauthor}
}
\makeatletter
\def\@maketitle{%
\newpage
\null
\vskip 1em%
\begin{center}%
\let \footnote \thanks
{\LARGE \@title \par}%
\vskip 1em%
%{\large \@date}%
\end{center}%
\par
\vskip 1em}
\makeatother

\usepackage{lipsum}  
\usepackage{cmbright}


%%%%%%%%%%%%%%%%%%%%%%%%%%%%%%%%%%%%%%%%%%%%%%%%%%%%%%%%%%%%%%%%%%%%%%%%%%%%%%%%%%%%%%%%%%%%%%%%%%%%%%%%%%%%%%%%%%%%%%%%%%%%%%%%%%%%%%%%%%%%%%%%%%%%%%%%%%%%%%%%%%%%%%%%%%%%%%%%%%%%%%%%%%%%%%%%%%%%%%

\begin{document}



\maketitle


%%%%%%%%%%%%%%%%%%% Datos del profesor, campus y aula

\begin{mybox}
\noindent\begin{tabular}{@{}ll}
	{\bfseries{Profesor}} & Dr. Jorge Velázquez Castro\\
	\textcolor{prussianblue}{\bfseries{Campus}} / \textcolor{trueblue}{\bfseries{Facultad}}  & \textcolor{prussianblue}{\bfseries{C.U. BUAP}} /  \textcolor{trueblue}{\bfseries{FCFM}} \\
	\textcolor{prussianblue}{\bfseries{Edificio}} /  \textcolor{trueblue}{\bfseries{Salón}}    & 	\textcolor{prussianblue}{\bfseries{EMA4}} /  \textcolor{trueblue}{\bfseries{407}}   
\end{tabular} 
\end{mybox} \vspace{1cm}

\tableofcontents

\vspace{2cm}

\section{Repaso de temas de microcanónico}


\begin{figure}[H]
    \centering
	    \includegraphics[width=0.7\textwidth]{1.png} % ajusta el ancho según convenga
    \caption{Descripción clara y concisa de la figura.}
    \label{fig:mi_figura2}
\end{figure}


\[  P_{c} (E) = \frac{1}{z}g(E) e^{-\beta E} \]

\begin{figure}[H]
    \centering
	    \includegraphics[width=0.7\textwidth]{2.png} % ajusta el ancho según convenga
    \caption{Descripción clara y concisa de la figura.}
    \label{fig:mi_figura2}
\end{figure}

Donde g(E) se le conoce como la \textbf{densidad de estados} y es equivalente a $\Omega{E}$.


\textbf{Equivalencia de ensambles}. En termodinámica no existe preferencia sobre algún sistema termodinámico. 

El sistema intercambia energía bajo un reservorio y todo lo construimos bajo esa hipótesis (estadística).


\begin{enumerate}
\item Primero encontremos el máximo de la grafica anterior
\end{enumerate}


Tenemos: 

Encontraremos el valor de $E^{\ast}$ que hace máxima a la probabilidad.

\[  \frac{\partial }{\partial E} \ln{P_{c}} (E^{\ast}) =0 \]

\[ \boxed{ \color{Cinnabar} \frac{\partial }{\partial E} \ln \left[  \frac{1}{z} g(E) e^{-\beta E}  \right]_{E^{\ast}} } =0\]



Tenemos:


\[ \frac{\partial }{\partial E} \left[ \ln{g(E)} - \beta{E} - \ln{Z}\right]_{E^{\ast}}  =0\]


Tenemos:


\[ \frac{ \partial  }{\partial E} \left[ \frac{ S(E) }{k_{B}} - \frac{1}{k_{B}T} \right]_{E^{\ast}} =0  \]


\[  \Rightarrow \quad \boxed{\color{Cinnabar} \frac{ \partial S }{\partial E}(E^{\ast}) = \frac{1}{T}  } \]

\textcolor{Cinnabar}{ \textbf{En un sistema en contacto térmico con el ambiente la representación entrópica trata con la energía más probable $E^{\ast}$}}

Anteriormente obtuvimos la expresión:


\[ U =- \frac{\partial }{\partial \beta} \ln{z} = \langle E \rangle \]

Tenemos:


Transformada de legendre de la entropía:


\[ \boxed{U =  - \frac{\partial}{\partial \frac{1}{T}} S\left[  \frac{1}{T} \right] = E } \]

EL resultado anterior en termodinámica nos daba la \textbf{energía}. Es equivalente a la representación del ensamble canónico. 


POr lo tanto tenemos la equivalencia siguiente:


\begin{equation}
\begin{split}
 \boxed{\color{Cinnabar} \frac{ \partial S }{\partial E}(E^{\ast}) = \frac{1}{T}  } &\equiv \boxed{U =  - \frac{\partial}{\partial \frac{1}{T}} S\left[  \frac{1}{T} \right] = E }
\end{split}
\end{equation}

\subsection{Tamaño de las fluctuaciones}


La desviación estándar $\sqrt{\sigma^{2}}$ de $P_{c}$ la obtenemos de:


\[ \sigma^{2}  = \langle E^{2} \rangle - \langle E \rangle^{2}  \]

Y lo podemos calcular con:

\[  \frac{dU}{d\beta}  = -\frac{1}{z} \int\limits_{0}^{\infty} g(E)E^{2}e^{-\beta E}dE + \frac{1}{z^{2}} \left( \int\limits_{0}^{\infty} g(E) Ee^{-\beta E}  dE\right)^{2}  \]


Esto es igual a:



\[ =- \left(  \langle E^{2} \rangle - \langle E\rangle^{2}   \right) =- \sigma^{2}  \]

\textbf{Nota}:

\[ \langle E \rangle \equiv U = \int\limits \frac{ E g(E)e^{-\beta E} }{z(\beta)}dE  \]


También:


\[  \beta = \frac{1}{k_{B} T}  \]



entonces:


\[ \sigma^{2} = - \frac{ \partial U }{\partial \beta} = K_{B} T^{2}\frac{dU}{dT}  \]

\[ \sigma^{2} = K_{B} T^{2}C_{v} \]

El ancho relativo de la distribución es:


\[   \frac{\sigma }{ \langle E\rangle } = \frac{ \sqrt{ \sigma^{2} } }{\langle E \rangle} = \frac{1}{U} \sqrt{K_{B}T^{2}C_{v} } \approx \frac{ \sqrt{N} }{N} = \frac{1}{\sqrt{N}} \]

Ya que:


\[ C_{v} \approx N, \quad U\approx N \]

Donde: $N$ es el número de partículas.

SI el aparato de medición es más preciso que $N$, entonces es importante considerar las \textbf{fluctuaciones}. Ya no podemos usar termodinámica clásica.


\textbf{Las fluctuaciones} decaen  en el compartamiento de :


\[ \frac{1}{N} \]

\begin{figure}[H]
    \centering
	    \includegraphics[width=0.7\textwidth]{3.png} % ajusta el ancho según convenga
    \caption{Descripción clara y concisa de la figura.}
    \label{fig:mi_figura2}
\end{figure}

La figura anterior nos dice cómo va cambiando el ancho con respecto a los valores del dominio, es decir es la tasa del cociente de:


\[ \frac{\sigma}{\langle E \rangle} \]

El cociente anterior es un resultado muy conocido, podemos saber cuando tomar en cuenta las fluctuaciones y cuando no tomarlas. Es algo parecido a por que la longitud térmica es conocida. Cuando la longitud térmica es del recipiente (caja) que contiene  a la partícula, entonces suceden efectos cuánticos.

Podemos tamnbién escribir $P_{c}$ en términos de cantidades termodinámicas.

Vamos a hacer una expansión (en serie de \textbf{Taylor}) alrededor del máximo de $P_{c}$ en $E^{\ast}$

Tenemos:


\[ \ln \left[  ge^{-\beta E}  \right] \approx \ln\left[ ge^{-\beta E}   \right]_{E^{\ast}} + \frac{1}{2} \frac{\partial^{2}}{\partial E^{2}} \ln\left[ ge^{-\beta E}    \right]_{E^{\ast}} \left(  E-E^{\ast} \right)^{2}   \]

\textit{Las \textbf{fluctuaciones} \underline{relativas} están directamente ralacionados con la mediciones de los aparatos de medición.}



Tenemos realizando la expansión que necesitabamos:


\[ \left[   \frac{1}{ K_{B} } S - \beta E \right]_{E^{\ast}} + \frac{1}{2} \frac{ \partial^{2} }{\partial E^{2}} \left[ \frac{1}{K_{B}}S  \right]_{E^{\ast}} \left(  E-E^{\ast}\right)^{2}  \]


Por otro lado:


\[ \frac{ \partial^{2}S }{\partial E^{2}} = \frac{\partial }{\partial E} \left( \frac{1}{T} \right) = \frac{\partial T}{\partial E} \frac{ d \frac{1}{T} }{dT} =- \frac{ \frac{1}{T^{2}} }{ \frac{ \partial E }{\partial T}   } =- \frac{1}{T^{2}} \frac{1}{C_{v}}      \]

Por lo que:

\begin{figure}[H]
    \centering
	    \includegraphics[width=0.7\textwidth]{4.png} % ajusta el ancho según convenga
    \caption{Descripción clara y concisa de la figura.}
    \label{fig:mi_figura2}
\end{figure}

\begin{figure}[H]
    \centering
	    \includegraphics[width=0.7\textwidth]{5.png} % ajusta el ancho según convenga
    \caption{Descripción clara y concisa de la figura.}
    \label{fig:mi_figura2}
\end{figure}


\section{Teorema del virial y teorema de equipartición}

Si $H(q_{s},p_{s})$ es la hamiltoniana, denotamos a $q_{s}$ y $p_{s}$ con $x_{i}$ donde $i=1, \ldots , 6N$

Entonces:


\[ \langle x_{i} \frac{dH}{dx_{k}} \rangle_{L} = \frac{1}{h^{3}N} \int\limits dx^{6N} \rho_{c}(\vec{x})x_{i} \frac{\partial H}{\partial x_{k}}  \]


Como $\rho_{c}(\vec{x})= \frac{1}{z}e^{-\beta H(\vec{x})}$



Y además:


\[ e^{-\beta H} \frac{\partial H }{\partial x_{k}} = - \frac{1}{\beta} \frac{ \partial }{\partial x_{k}} e^{-\beta H} \]


Entonces:

\begin{figure}[H]
    \centering
	    \includegraphics[width=0.7\textwidth]{6.png} % ajusta el ancho según convenga
    \caption{Descripción clara y concisa de la figura.}
    \label{fig:mi_figura2}
\end{figure}


Por asumsión física $H \to \infty$ en $x_{max}$ y $x_{min}$


Entonces:


\[  \langle  x_{i} \frac{\partial H }{\partial x_{k}}  \rangle = \frac{ \delta_{ik} }{\beta} = \delta{ik}K_{B}T   \]

Caso (a):

Si $x_{i}$ es una $q_{i}$, $\frac{\partial H}{\partial q_{i}} = -\dot{p}_{i}$


Es decir:


\[ \langle  x_{i} \frac{\partial H}{\partial x_{i}} \rangle = - \langle q_{i}\dot{p}_{i}\rangle = - \langle q_{i} \underbrace{F_{i}}_{\text{Fuerza generalizada}} \rangle = K_{B}T \]

\begin{figure}[H]
    \centering
	    \includegraphics[width=0.7\textwidth]{7.png} % ajusta el ancho según convenga
    \caption{Descripción clara y concisa de la figura.}
    \label{fig:mi_figura2}
\end{figure}

Por otro lado:


\[ \overbrace{T}^{\text{E. Cinética}} = \frac{ \vec{p}^{2} }{2m} = \frac{p_{x}^{2}}{2m} + \frac{p_{y}^{2}}{2m} + \frac{p^{2}_{z}}{2m} \]

\[\langle T \rangle  = \frac{3}{2} \langle \frac{p_{i}^{2}}{m} \rangle \]

Si:

\[ F_{i} = -\nabla{v_{i}}; \quad v_{i} = \text{potencial mecánico}   \]

Por ejemplo

\[ v=r^{\alpha} \]

\begin{figure}[H]
    \centering
	    \includegraphics[width=0.7\textwidth]{8.png} % ajusta el ancho según convenga
    \caption{Descripción clara y concisa de la figura.}
    \label{fig:mi_figura2}
\end{figure}

Este resultado es igual a:

\[ \alpha \langle v_{i} \rangle = - 3K_{B} T\]

La temperatura es proporcional a la energía potencial del sistema. Es una forma intuitiva de ver a la temperatura, pero también en la mayoría de los casos es proporcional al \textbf{potencial}.  

\subsection{Caso particular $\alpha=2$}



\begin{figure}[H]
    \centering
	    \includegraphics[width=0.7\textwidth]{9.png} % ajusta el ancho según convenga
    \caption{Descripción clara y concisa de la figura.}
    \label{fig:mi_figura2}
\end{figure}

\begin{figure}[H]
    \centering
	    \includegraphics[width=0.7\textwidth]{10.png} % ajusta el ancho según convenga
    \caption{Descripción clara y concisa de la figura.}
    \label{fig:mi_figura2}
\end{figure}




\end{document}